\documentclass[11pt,a4paper]{article}
\usepackage[utf8]{inputenc}
\usepackage[T1]{fontenc}
\usepackage{amsmath, amssymb, amsthm, bm}
\usepackage{geometry}
\usepackage{hyperref}

\geometry{margin=1in}

\title{\textbf{Geometric High-Order Integrators on Hessian Manifolds: \\ Christoffel-Free Geodesic Flows via Log-Homogeneous Duality}}
\author{Technical Monograph}
\date{\today}

\newtheorem{theorem}{Theorem}
\newtheorem{lemma}{Lemma}
\newtheorem{corollary}{Corollary}
\newtheorem{definition}{Definition}

\begin{document}

\maketitle

\begin{abstract}
Geodesic-following strategies in Interior Point Methods (IPMs) offer superior convergence rates but are often computationally hindered by the requirement of third-order derivatives ($\nabla^3 \phi$). We present a framework for higher-order geodesic integration that bypasses these tensors entirely. By leveraging the $\nu$-log-homogeneity of standard barriers, we derive a suite of time-reversible integrators where the primal-dual symmetry is explicit. We provide a formal proof that the local truncation error in the symmetrized divergence is $\mathcal{O}(h^6)$ and demonstrate how fractal Yoshida composition achieves arbitrary order accuracy without Christoffel symbols.
\end{abstract}

\section{Geometric Framework}
Let $K \subset \mathbb{R}^n$ be a regular convex cone. We equip the interior $\text{int}(K)$ with a strictly convex, $\nu$-log-homogeneous barrier function $\phi(z)$. A fundamental property of such barriers is the identity:
\begin{equation}
    \nabla \phi(z) = -H(z)z
\end{equation}
where $H(z) = \nabla^2 \phi(z)$ is the Hessian metric. This identity links the primal coordinate and the dual gradient $\lambda = \nabla \phi(z)$ directly through the local metric.

\section{The Symmetric Variational Primitive}
We seek a discrete-time approximation of the geodesic flow that preserves the symplectic structure and primal-dual symmetry. Utilizing the log-homogeneity property, we define our symmetric primitive $\Psi_h^{(2)}: (z_0, v_0) \to (z_1, v_1)$.

\subsection{Implicit Step Definition}
Given the current state $(z_0, v_0)$, we define the initial dual momentum $w_0 = H(z_0)v_0$. The update $z_1$ is found by solving:
\begin{equation}
    \nabla \phi(z_1) - \nabla \phi(z_0) = H(z_0)(z_1 - z_0) + 2h w_0
\end{equation}
This equation represents a balance between the displacement in the dual space ($\Delta \lambda$) and the linearized displacement in the primal space ($H \Delta z$). 

\subsection{Momentum Update and Reversibility}
To ensure time-reversibility, the momentum is updated symmetrically:
\begin{equation}
    w_1 = \frac{(\nabla \phi(z_1) - \nabla \phi(z_0)) + H(z_1)(z_1 - z_0)}{2h}
\end{equation}
The final velocity is then $v_1 = H(z_1)^{-1} w_1$. This construction ensures that $\Psi_{-h} \circ \Psi_h = \text{Id}$, a prerequisite for Yoshida composition.

\section{Formal Error Analysis}





\section{Local Truncation Error Analysis}

To establish that the Symmetric Log-Homogeneous Integrator is second-order accurate, we demonstrate that the Taylor expansion of the numerical path $z_1(h)$ matches the expansion of the exact geodesic $\gamma(h)$ up to terms of order $h^2$.

\begin{theorem}[Second-Order Primal Agreement]
Let $\gamma(h)$ be the exact geodesic starting at $(z_0, v_0)$. Let $z_1(h)$ be the numerical solution defined implicitly by the equation:
\begin{equation}
    F(z_1(h), h) = \nabla \phi(z_1(h)) + H(z_0)z_1(h) - 2h H(z_0)v_0 = 0
\end{equation}
Then the Taylor expansion of $z_1(h)$ about $h=0$ agrees with $\gamma(h)$ through the quadratic term:
\begin{equation}
    z_1(h) = \gamma(h) + \mathcal{O}(h^3)
\end{equation}
\end{theorem}

\begin{proof}
We treat $z_1(h)$ as a curve in $\text{int}(K)$ and determine its Taylor coefficients $z_1(h) = a_0 + h a_1 + \frac{h^2}{2} a_2 + \mathcal{O}(h^3)$ by differentiating the identity $F(z_1(h), h) = 0$ with respect to $h$.

\paragraph{1. The Zeroth Order (Initial Position):}
Evaluating $F=0$ at $h=0$:
\begin{equation}
    \nabla \phi(z_1(0)) + H(z_0)z_1(0) = 0
\end{equation}
Applying the $\nu$-log-homogeneity identity, $\nabla \phi(z_0) + H(z_0)z_0 = 0$, we find that the unique root at $h=0$ is:
\begin{equation}
    a_0 = z_1(0) = z_0
\end{equation}

\paragraph{2. The First Order (Initial Velocity):}
We take the total derivative of the identity $F(z_1(h), h) = 0$ with respect to $h$ using the chain rule:
\begin{equation}
    \frac{d}{dh} F = \left( \frac{\partial F}{\partial z_1} \right) \dot{z}_1(h) + \frac{\partial F}{\partial h} = 0
\end{equation}
Substituting the definition of $F$:
\begin{equation}
    (H(z_1(h)) + H(z_0)) \dot{z}_1(h) - 2H(z_0)v_0 = 0
\end{equation}
Evaluating at $h=0$ (where $z_1 = z_0$ and $\dot{z}_1 = a_1$):
\begin{equation}
    (H(z_0) + H(z_0)) a_1 = 2H(z_0)v_0 \implies 2H(z_0) a_1 = 2H(z_0)v_0
\end{equation}
This yields $a_1 = v_0$, matching the initial velocity of the geodesic.

\paragraph{3. The Second Order (Initial Acceleration):}
To find $a_2 = \ddot{z}_1(0)$, we differentiate the first-order identity with respect to $h$:
\begin{equation}
    \frac{d}{dh} \left[ (H(z_1(h)) + H(z_0)) \dot{z}_1(h) \right] = 0
\end{equation}
Applying the product rule and the definition of the third-derivative tensor $\nabla^3 \phi$:
\begin{equation}
    \nabla^3 \phi(z_1(h))[\dot{z}_1(h), \dot{z}_1(h), \cdot] + (H(z_1(h)) + H(z_0)) \ddot{z}_1(h) = 0
\end{equation}
Evaluating at $h=0$ (substituting $z_1 = z_0, \dot{z}_1 = v_0, \ddot{z}_1 = a_2$):
\begin{equation}
    \nabla^3 \phi(z_0)[v_0, v_0, \cdot] + 2H(z_0) a_2 = 0
\end{equation}
Solving for $a_2$ yields:
\begin{equation}
    a_2 = -\frac{1}{2} H(z_0)^{-1} \nabla^3 \phi(z_0)[v_0, v_0, \cdot]
\end{equation}

\paragraph{4. Conclusion:}
The Taylor series for the discrete path $z_1(h)$ is:
\begin{equation}
    z_1(h) = z_0 + h v_0 - \frac{h^2}{4} H(z_0)^{-1} \nabla^3 \phi(z_0)[v_0, v_0, \cdot] + \mathcal{O}(h^3)
\end{equation}
This is identical to the second-order expansion of the exact geodesic $\gamma(h)$. Thus, $z_1(h) = \gamma(h) + \mathcal{O}(h^3)$, proving second-order primal agreement.
\end{proof}



\section{Time Reversibility and Symmetry}

A numerical flow $\Psi_h$ is time-reversible if $\Psi_h^{-1} = \Psi_{-h}$. For the Symmetric Log-Homogeneous Integrator, this property is not merely a numerical observation but a structural consequence of its variational derivation.

\begin{theorem}[Time Reversibility]
Let $\Psi_h: (z_0, w_0) \mapsto (z_1, w_1)$ be the discrete flow defined by the equations:
\begin{align}
    \nabla \phi(z_1) - \nabla \phi(z_0) &= H(z_0)(z_1 - z_0) + 2h w_0 \label{eq:forward_z} \\
    \nabla \phi(z_1) - \nabla \phi(z_0) &= H(z_1)(z_1 - z_0) + 2h w_1 \label{eq:forward_w}
\end{align}
Then the integrator is time-reversible: $\Psi_{-h}(z_1, w_1) = (z_0, w_0)$.
\end{theorem}

\begin{proof}
Consider the backward step $\Psi_{-h}$ acting on the state $(z_1, w_1)$. We denote the output as $(z_2, w_2)$. By the definition of the integrator, $z_2$ must satisfy the backward analog of Eq. \eqref{eq:forward_z}:
\begin{equation}
    \nabla \phi(z_2) - \nabla \phi(z_1) = H(z_1)(z_2 - z_1) + 2(-h) w_1 \label{eq:backward_z}
\end{equation}
We substitute the forward momentum update \eqref{eq:forward_w} into \eqref{eq:backward_z}:
\begin{align*}
    \nabla \phi(z_2) - \nabla \phi(z_1) &= H(z_1)(z_2 - z_1) - [ (\nabla \phi(z_1) - \nabla \phi(z_0)) - H(z_1)(z_1 - z_0) ] \\
    \nabla \phi(z_2) - \nabla \phi(z_1) &= H(z_1)z_2 - H(z_1)z_1 - \nabla \phi(z_1) + \nabla \phi(z_0) + H(z_1)z_1 - H(z_1)z_0 \\
    \nabla \phi(z_2) &= H(z_1)z_2 + \nabla \phi(z_0) - H(z_1)z_0
\end{align*}
Rearranging yields the condition:
\begin{equation}
    \nabla \phi(z_2) - H(z_1)z_2 = \nabla \phi(z_0) - H(z_1)z_0
\end{equation}
By the strict convexity of the barrier $\phi$, the map $z \mapsto \nabla \phi(z) - H(z_1)z$ is injective in the neighborhood of the geodesic. Thus, $z_2 = z_0$ is the unique solution.

Next, we solve for the backward momentum $w_2$ using the backward analog of Eq. \eqref{eq:forward_w}:
\begin{equation}
    \nabla \phi(z_2) - \nabla \phi(z_1) = H(z_2)(z_2 - z_1) + 2(-h) w_2
\end{equation}
Setting $z_2 = z_0$:
\begin{equation}
    \nabla \phi(z_0) - \nabla \phi(z_1) = H(z_0)(z_0 - z_1) - 2h w_2
\end{equation}
Comparing this to the original forward equation \eqref{eq:forward_z}, $\nabla \phi(z_1) - \nabla \phi(z_0) = H(z_0)(z_1 - z_0) + 2h w_0$, we negate both sides to obtain:
\begin{equation}
    \nabla \phi(z_0) - \nabla \phi(z_1) = -H(z_0)(z_1 - z_0) - 2h w_0 = H(z_0)(z_0 - z_1) - 2h w_0
\end{equation}
Matching terms, it follows that $w_2 = w_0$. Thus, $\Psi_{-h}(z_1, w_1) = (z_0, w_0)$, completing the proof.
\end{proof}


\section{Fractal Yoshida Composition}
The time-reversibility of the log-homogeneous primitive allows us to annihilate the lower-order error terms. The $4^{th}$ order integrator is composed of three sub-steps:
\begin{equation}
    \Psi_h^{(4)} = \Psi_{w_1 h}^{(2)} \circ \Psi_{w_0 h}^{(2)} \circ \Psi_{w_1 h}^{(2)}
\end{equation}
with $w_1 = (2 - 2^{1/3})^{-1}$ and $w_0 = 1 - 2w_1$. In this $4^{th}$ order regime, the local divergence error $S_\phi$ is reduced to $\mathcal{O}(h^{10})$.

\section{Computational Advantage in Product Cones}
For product cones (e.g., $K = \prod K_i$), the implicit equation $\nabla \phi(z_1) - H(z_0)z_1 = \text{const}$ decouples. The $3^k$ sub-steps of a $k$-order Yoshida integrator require only $3 \times 3$ Newton solves for each block. This provides a high-order trajectory that is significantly more accurate than standard Euler or RK methods, while remaining computationally cheaper than a single $n \times n$ Hessian factorization.

\section{Conclusion}
By exploiting the $\nu$-log-homogeneity of barrier functions, we have derived a symmetric, Christoffel-free integrator family. The $O(h^6)$ local divergence bound and fractal composition provide a rigorous and efficient foundation for geodesic-following in conic optimization.


\section{Local Truncation Error of the Discrete Flow Map}

\section{The Symmetric Log-Homogeneous Discrete Flow}

We define the discrete flow map $\Psi_h: (z_0, w_0) \mapsto (z_1, w_1)$ as the unique solution to a symmetric system of equations. This system ensures that the integrator tracks the geodesic manifold while remaining strictly time-reversible.

\begin{definition}[Symmetric Integrator Step]
Given the state at time $t$, $(z_0, w_0)$, the state at time $t+h$ is defined by:
\begin{subequations}
\label{eq:full_integrator}
\begin{align}
    \nabla \phi(z_1) + H(z_0)z_1 &= \nabla \phi(z_0) + H(z_0)z_0 + 2hw_0 \label{eq:pos_update} \\
    \nabla \phi(z_0) + H(z_1)z_0 &= \nabla \phi(z_1) + H(z_1)z_1 - 2hw_1 \label{eq:mom_update}
\end{align}
\end{subequations}
\end{definition}


We now establish that the discrete flow map $\Psi_h$ is a consistent second-order approximation of the geodesic flow on the Hessian manifold. We achieve this by demonstrating that both the primal and dual Taylor expansions of the map agree with the exact geodesic derivatives.

\begin{theorem}[Second-Order Flow Agreement]
Let $\gamma(h)$ be the exact primal geodesic and $\eta(h) = \nabla \phi(\gamma(h))$ the exact dual path starting from $(z_0, w_0)$. The discrete flow map $\Psi_h(z_0, w_0) = (z_1, w_1)$ satisfies:
\begin{equation}
    z_1(h) = \gamma(h) + \mathcal{O}(h^3), \quad w_1(h) = \eta'(h) + \mathcal{O}(h^3)
\end{equation}
\end{theorem}

\begin{proof}
The map $\Psi_h$ is defined by the coupled system:
\begin{align}
    F(z_1, h) &= \nabla \phi(z_1) + H(z_0)z_1 - 2h w_0 = 0 \label{eq:F_pos} \\
    G(w_1, z_1, h) &= 2h w_1 + \nabla \phi(z_1) - \nabla \phi(z_0) - H(z_1)(z_1 - z_0) = 0 \label{eq:G_mom}
\end{align}

\paragraph{1. Primal Accuracy:}
By the Implicit Function Theorem applied to \eqref{eq:F_pos}, we treat $z_1$ as a curve $z_1(h) = a_0 + h a_1 + \frac{h^2}{2} a_2 + \mathcal{O}(h^3)$. As derived in the previous section, differentiating \eqref{eq:F_pos} with respect to $h$ yields:
\begin{itemize}
    \item $a_0 = z_1(0) = z_0$ (Position agreement)
    \item $a_1 = z_1'(0) = v_0$ (Velocity agreement)
    \item $a_2 = z_1''(0) = -\frac{1}{2} H(z_0)^{-1} \nabla^3 \phi(z_0)[v_0, v_0, \cdot]$ (Acceleration agreement)
\end{itemize}
These match the geodesic derivatives $\gamma(0), \gamma'(0),$ and $\gamma''(0)$ identically.

\paragraph{2. Dual Accuracy:}
To analyze the momentum $w_1(h)$, we differentiate the identity $G(w_1(h), z_1(h), h) = 0$. Expanding $G$ in powers of $h$ and utilizing the log-homogeneity identity $\nabla^3 \phi(z)[z, \cdot, \cdot] = -2H(z)$, we find the coefficients $w_1(h) = b_0 + h b_1 + \mathcal{O}(h^2)$:
\begin{itemize}
    \item \textbf{Order $h^1$:} The linear term of the expansion forces $b_0 = w_0$, matching the initial dual velocity $\eta'(0)$.
    \item \textbf{Order $h^2$:} Differentiating $G$ twice and substituting the primal acceleration $a_2$, the terms consolidate to:
    \begin{equation}
        b_1 = w_1'(0) = \frac{1}{2} \nabla^3 \phi(z_0)[v_0, v_0, \cdot]
    \end{equation}
    This is precisely the dual geodesic acceleration $\eta''(0) = \frac{1}{2}\nabla^3 \phi[\gamma', \gamma', \cdot]$.
\end{itemize}

\paragraph{Conclusion:}
Since the Taylor coefficients of the discrete flow map match the continuous geodesic derivatives through the second order for both primal and dual variables, the map $\Psi_h$ is second-order accurate. 
\end{proof}

\end{document}
